For matrices $A \in \FF^{n\times m}$ and $B \in \FF^{n \times k}$ we denote by $[A; B] \in \FF^{n \times (m+k)}$ the concatenation of $A$ and $B$. For any matrix $A$ we denote by $A|_{[d]}$ the matrix consisting of the first $d$ columns of $A$. We denote by $\col(A)$ the span of the columns of $A$.

%\ILnote{Comment 9. Fixed below the instances of $[A;(BU)|_{[s]}$ which all missed the closing square bracket (missed it, and not placed the restriction inside, as the comment says).}
\begin{lemma}\label{lem:colspan}
    Let $A \in \FF^{n\times m}$ and $B \in \FF^{n \times k}$. There is a nonzero %$k^2$-variable 
    polynomial $p\in \FF[x_{1,1},x_{1,2},,\ldots,x_{i,j},\dots,x_{k,k}]$ of degree $s = \rank([A;B]) - \rank(A)$ such that for any matrix $U \in \FF^{k\times k}$ satisfying $p(U)\neq 0$, we have that
    \[
    \col([A;B]) = \col([A; (BU)|_{[s]}]).
    \]
\end{lemma}
\begin{proof}
    Let $r = \rank([A;B])$. Then, there is an $r\times r$ submatrix of $[A;B]$ with rank~$r$ induced by $r-s$ columns of~$A$ and~$s$ columns of~$B$.
    Let~$R\subseteq [n]$ be the set of row indices of this submatrix, and let $C_A\subseteq[m], C_B\subseteq \{m+1,\dots,m+k\}$ be the respective column indices of~$A$ and~$B$ in this submatrix.
    Let $p(X)$ be the determinant of the submatrix of $[A;BX]$ induced by the rows in~$R$ and columns in $C_A\cup \{m+1,\dots,m+s\}$, where $X = (x_{i,j})_{i,j=1}^k$ is a matrix of variables.
    
    We claim that~$p$ is nonzero. There is a permutation matrix~$V\in \FF^{k \times k}$ that swaps the columns in~$C_B$ with the columns in $\{m+1,\dots, m+s\}$. Since the submatrix of $[A;BV]$ induced by the rows in~$R$ and columns in $C_A\cup \{m+1,\dots,m+s\}$ is then of full rank, it follows that~$p(V)\ne 0$. 
        
    Let $U \in \FF^{k \times k}$ be a matrix such that $p(U) \neq 0$. 
    Then we have $\rank([A; (BU)|_{[s]}]) \geq r$. From $\col([A; (BU)|_{[s]}])\subseteq \col([A;B])$ and $\rank([A;B]) = r$, we find the equality $\col([A;B]) = \col([A; (BU)|_{[s]}])$.
\end{proof}

% \refnote{10}{Wouldn't it be more intuitive to merge 2.4-2.5 into one lemma where you directly add $s_i$ columns at a time?}

% \JBnote{The proofs were merged and only the second lemma remains, as suggested.}

% \begin{lemma}\label{lem:schwartz-zippel}
%     Let $A_1, \ldots, A_\ell \in \FF^{n \times k}$ be matrices. Suppose $|\FF| > \rank([A_1; \cdots; A_\ell])$. Let 
%     \[
%     s_i = \rank([A_1; \cdots ; A_{i}]) - \rank([A_1; \cdots ; A_{i-1}]).
%     \]
%     There is a matrix $U \in \FF^{k \times k}$ such that for all $i \in [\ell]$ we have 
%     \[
%     \col([A_1 ; A_2 ; \cdots ; A_{i-1} ; (A_i U)|_{[s_i]}]) = \col([A_1; \cdots ; A_i]).
%     \]
% \end{lemma}
% \begin{proof}
% From Lemma~\ref{lem:colspan}, we get that for each $i\in [\ell]$, there exists a nonzero~$k^2$-variable polynomial $p_i(X)\in \FF[x_{1,1},x_{1,2},\dots,x_{k,k}]$ of degree~$s_i$ such that for any $U\in \FF^{k\times k}$ satisfying~$p_i(U)\ne 0$,
% \[
%     \col([A_1 ; A_2 ; \cdots ; A_{i-1} ; (A_i U)|_{[s_i]}]) = \col([A_1; \cdots ; A_i]).
%     \]
% Let $q(X)=  p_1(X)\cdots p_\ell(X)$.
% This is a nonzero polynomial of degree $s_1 + \cdots +s_\ell = \rank([A_1; \cdots; A_\ell])<|\FF|$.
% It follows from the Schwartz--Zippel lemma \cite{DEMILLO1978193, schwartz1980fast, zippel1979probabilistic} that there is a~$U\in \FF^{k\times k}$ such that~$q(U) \ne 0$.
% This is to say that for all~$i\in[\ell]$ simultaneously,~$p_i(U)\ne 0$, which implies the result.
% \end{proof}

\begin{lemma}\label{lem:max-rank-induction}
    Let $A_1, \ldots, A_m \in \FF^{n \times k}$ be matrices. Suppose $|\FF| > \rank([A_1; \cdots; A_m])$. Let 
    \[
    s_i = \rank([A_1; \cdots ; A_{i}]) - \rank([A_1; \cdots ; A_{i-1}]).
    \]
    There is a matrix $U \in \FF^{k \times k}$ such that
\begin{equation}\label{eq:cspan2}
    \col([(A_1 U)|_{[s_1]} ; (A_2 U)|_{[s_2]} ; \cdots ; (A_m U)|_{[s_m]}]) = \col([A_1; \cdots ; A_m]).
\end{equation}
\end{lemma}
% \begin{proof}
% %
% Let~$U\in \FF^{k\times k}$ be a matrix as in Lemma~\ref{lem:schwartz-zippel}. We give a proof by induction.
% The base case is $\col((A_1 U)|_{[s_1]}) = \col(A_1)$,
% which follows directly from \autoref{lem:schwartz-zippel}.
% For the induction step, we assume
% \begin{equation}\label{eq:ih}
%     \col([A_1 ; A_2 ; \cdots ; A_i]) = \col([(A_1U)|_{[s_1]} ; \cdots ; (A_iU)|_{[s_i]}]).
% \end{equation}
% We now use that if $\col(A) = \col(B)$, then $\col([A ; C]) = \col([B;C])$ holds for any matrices $A,B,C$ with the same number of rows. Applying this to \eqref{eq:ih} gives
% \begin{multline}\label{eq:x}
%     \col([A_1 ; A_2 ; \cdots ; A_i; (A_{i+1} U)|_{[s_{i+1}]}])\\ = \col([(A_1U)|_{[s_1]} ; \cdots ; (A_iU)|_{[s_i]}; (A_{i+1} U)|_{[s_{i+1}]}]).
% \end{multline}
% From \autoref{lem:schwartz-zippel} we know that 
% \begin{equation}\label{eq:sz}
%     \col([A_1; \cdots ; A_{i+1}]) = \col([A_1 ; A_2 ; \cdots ; A_i ; (A_{i+1} U)|_{[s_{i+1}]}]).
% \end{equation}
% %
% %
% Then \eqref{eq:x} and \eqref{eq:sz} together give
%     \[
%     \col([A_1 ; A_2 ; \cdots ; A_i; A_{i+1}]) = \col([(A_1U)|_{[s_1]} ; \cdots ; (A_iU)|_{[s_i]}; (A_{i+1}U)|_{[s_{i+1}]}]).
%     \]
% This proves the claim by induction.
% \end{proof}


\begin{proof}
We begin by showing that there exists a matrix $U \in \FF^{k \times k}$ such that for all $i \in [m]$ we have 
\begin{equation}\label{eq:cspan1}
    \col([A_1 ; A_2 ; \cdots ; A_{i-1} ; (A_i U)|_{[s_i]}]) = \col([A_1; \cdots ; A_i]).
\end{equation}
From Lemma~\ref{lem:colspan}, we get that for each $i\in [m]$, there exists a nonzero~$k^2$-variable polynomial $p_i(X)\in \FF[x_{1,1},x_{1,2},\dots,x_{k,k}]$ of degree~$s_i$ such that for any $U\in \FF^{k\times k}$ satisfying~$p_i(U)\ne 0$,
\[
    \col([A_1 ; A_2 ; \cdots ; A_{i-1} ; (A_i U)|_{[s_i]}]) = \col([A_1; \cdots ; A_i]).
    \]
Let $q(X)=  p_1(X)\cdots p_m(X)$.
This is a nonzero polynomial of degree $s_1 + \cdots +s_m = \rank([A_1; \cdots; A_m])<|\FF|$.
It follows from the Schwartz--Zippel lemma \cite{DEMILLO1978193, schwartz1980fast, zippel1979probabilistic} that there is a~$U\in \FF^{k\times k}$ such that~$q(U) \ne 0$.
This is to say that for all~$i\in[m]$ simultaneously,~$p_i(U)\ne 0$, which implies the existence of the desired matrix~$U$.

Let~$U\in \FF^{k\times k}$ be a matrix so that~\eqref{eq:cspan1} holds for each $i\in [m]$.
Next, we show by induction on~$i\in[m]$ that~$U$ also satisfies~\eqref{eq:cspan2}.
%Let~$U\in \FF^{k\times k}$ be a matrix as in Lemma~\ref{lem:schwartz-zippel}. We give a proof by induction.
The base case is $\col((A_1 U)|_{[s_1]}) = \col(A_1)$,
which follows directly.
For the induction step, we assume
\begin{equation}\label{eq:ih}
    \col([A_1 ; A_2 ; \cdots ; A_i]) = \col([(A_1U)|_{[s_1]} ; \cdots ; (A_iU)|_{[s_i]}]).
\end{equation}
We now use that if $\col(A) = \col(B)$, then $\col([A ; C]) = \col([B;C])$ holds for any matrices $A,B,C$ with the same number of rows. Applying this to \eqref{eq:ih} gives
\begin{equation}\label{eq:x}
    \col([A_1 ; A_2 ; \cdots ; A_i; (A_{i+1} U)|_{[s_{i+1}]}]) = \col([(A_1U)|_{[s_1]} ; \cdots ; (A_iU)|_{[s_i]}; (A_{i+1} U)|_{[s_{i+1}]}]).
\end{equation}
From \eqref{eq:cspan1} we know that 
\begin{equation}\label{eq:sz}
    \col([A_1; \cdots ; A_{i+1}]) = \col([A_1 ; A_2 ; \cdots ; A_i ; (A_{i+1} U)|_{[s_{i+1}]}]).
\end{equation}
%
%
Then \eqref{eq:x} and \eqref{eq:sz} together give
    \[
    \col([A_1 ; A_2 ; \cdots ; A_i; A_{i+1}]) = \col([(A_1U)|_{[s_1]} ; \cdots ; (A_iU)|_{[s_i]}; (A_{i+1}U)|_{[s_{i+1}]}]).
    \]
This proves the claim by induction.
\end{proof}



\begin{proof}[Proof of \autoref{th:uncertainty-principle-n}]
We will prove $n_1 \leq \subrank_2(T) \subrank_3(T)$.
Let $A_1, \ldots, A_{n_3} \in \FF^{n_1 \times n_2}$ be the $3$-slices of $T$.
Let~$U$ be as in \autoref{lem:max-rank-induction}.
Let 
\[
s_i = \rank([A_1; \cdots ; A_{i}]) - \rank([A_1; \cdots ; A_{i-1}]).
\]
Conciseness of~$T$ gives that $s_1 + \cdots + s_{n_3} = n_1$.
It follows from this and \autoref{lem:max-rank-induction} that the columns of the submatrices $(A_iU)|_{[s_i]}$ are linearly independent.
This implies that for each $i\in[n_3]$, 
there is a 3-slice of rank at least $s_i$,
%
and in particular $\subrank_3(T)\geq\max_i s_i$. %
%
Let $S$ be the tensor with 3-slices $A_1U, A_2U, \ldots, A_{n_3}U$. Then $\subrank_2(S) \leq \subrank_2(T)$ since the 2-slices of $S$ are linear combinations of the 2-slices of $T$. 
The matrix $M = [(A_1U)|_{[1]};\cdots ; (A_{n_3}U)|_{[1]}]$ consisting of every first column of the 3-slices of $S$ equals the first 2-slice of~$S$. Since the linearly independent columns in the 3-slices of~$S$ are flushed to the left, we find that $M$ has at least $|\{i\in [n_3]: s_i\ne 0\}|$ linearly independent columns, and thus $\rank(M) \geq |\{i\in [n_3]: s_i\ne 0\}|$,
%
%
which gives a lower bound on~$\subrank_2(S)$ and thus on $\subrank_2(T)$.
Since
\[
n_1  = s_1 + \cdots + s_{n_3}\leq |\{i\in [n_3]: s_i\ne 0\}|\, \max_i s_i \leq \subrank_2(T) \subrank_3(T),
\]
the result follows. %\ILnote{I think we should add: "for $\ell_1=2,\ell_2=3,\ell_3=1$ and similarly for other choices." and add to the last equation above the inequality $\leq\subrank_2(T)\subrank_3(T)$}.
%
%
%
%
%
%
%
%
%
%
%
\end{proof}
